\subsection{Эмпирическая функция распределения. Теорема Гливенко.}
Пусть задано некоторое вероятностное пространство $(\Omega, \mathcal{F}, \P)$ и последовательность $i.i.d.$ случайных величин $\{\xi_j\}_{j = 1}^\infty$ на нём.
\begin{definition}
    Определим последовательность эмпирических функций распределения следующим образом:
    \[
        \forall n \in \N\colon F_n(x) = \dfrac{1}{n}\sum\limits_{j = 1}^n \Ind_{\xi_j < x}.
    \]
\end{definition}
\begin{note}
    Для эмпирических функций распределения можно заметить следующие свойства:
    \begin{itemize}
        \item $F_n$ обладает всеми свойствами функции распределения;
        \item $\Ex F_n(x) = \frac{1}{n}\Ex I_{\xi_j < x} = \P(\xi < x) = F_\xi(x)$;
        \item Согласно второй теореме Колмогорова последовательность эмпирических функций распределения сходится к $F_\xi$:
        \[
            \forall x \hookrightarrow F_n(x) \xrightarrow{a.s.} F_\xi(x), n \ra \infty.
        \]
        \item Если $F_\xi$ -- непрерывная функция распределения, то сходимость в пункте выше является равномерной.
    \end{itemize}
\end{note}
\begin{theorem}[Гливенко, основная теорема математической статистики]
    Для последовательности эмпирических функций распределения справедлива следующая сходимость:
    \[
        d_n = \sup\limits_{x \in \R} |F_n(x) - F_\xi(x)| \xrightarrow{a.s.} 0, n \ra \infty.
    \]
\end{theorem}
\begin{note}
    В дальнейшем доказательстве мы считаем $F_\xi(x)$ -- непрерывной функцией распределения.
\end{note}
\begin{proof}
    Заметим, что в силу плотности $\Q$ в $\R$ мы можем определить $d_n$ следующим образом:
    \[
        d_n = \sup\limits_{x \in \Q} |F_n(x) - F_\xi(x)|,
    \]
    а значит $d_n$ -- случайная величина (как супремум последовательности случайных величин). \\
    Определим множество $\mathcal{L}_k$ следующим образом:
    \[
        \mathcal{L}_k = \{w \mid F_n(x_k) \ra F_\xi(x_k), n \ra \infty\},
    \]
    где $x_k \in \Q$. \\
    Как ранее было замечено, по второй теореме Колмогорова имеет место сходимость 
    \[
        \forall x \hookrightarrow F_n(x) \xrightarrow{a.s.} F_\xi(x), n \ra \infty,
    \]
    в том числе и при $x = x_k$, а значит $\forall x_k \in \Q \hookrightarrow\P(\mathcal{L}_k) = 1 $. \\
    Теперь, если мы рассмотрим множество $\mathcal{L} = \bigcap\limits_{k \in \N} \mathcal{L}_k$, то можно заметить, что его мера равна единице. 
    Поскольку 
    \[
        \P(\mathcal{L}) = \P\biggr(\bigcap\limits_{x_k \in \Q}\{w \mid F_n(x) \rightarrow F_\xi(x), n \ra \infty\}\biggr) = 1,
    \]
    то и 
    \[
        \P\biggr(\bigcap\limits_{x_k \in \Q} \{w \mid \sup\limits_{x \in \Q}|F_n(x) - F_\xi(x)| \ra 0, n \ra \infty\}\biggr) = 1
    \]
    что и требовалось доказать.
\end{proof}